\section{The Sources and Their Weighing}

Everything in this study depends on being clear, before the life begins,
about what each class of evidence is and what it can carry.

\subsection{The seven undisputed letters}

The bedrock is the set of seven letters whose direct Pauline authorship
commands effectively universal assent, ancient and modern: Romans,
1~Corinthians, 2~Corinthians, Galatians, Philippians, 1~Thessalonians, and
Philemon. They were written within roughly a decade---on the chronology
defended below, from about 50 to the late 50s or early 60s---to living
communities and one householder, in the middle of disputes whose other side
survives only as Paul reports it. They are class-A evidence in this
series' scheme: subject-authored, and therefore uniquely direct; but
``authored'' is not ``neutral.'' Each letter is occasional (it answers a
situation, not a biographer's questionnaire), rhetorical (it argues, pleads,
shames, and defends), and selective (Paul tells the Galatians about his past
precisely because his independence from Jerusalem is under attack, Gal
1:11--12). The autobiographical passages---Galatians 1--2, Philippians
3:4--6, 1~Corinthians 15:8--10, 2~Corinthians 11--12---are the most precious
paragraphs in the record, and every one of them is fighting for something.
That does not make them false; it makes them testimony, to be
cross-examined like any testimony.

\subsection{The disputed letters}

Six further letters bear Paul's name in the canon: 2~Thessalonians,
Colossians, Ephesians, and the three Pastorals (1--2~Timothy, Titus). Their
canonical standing is not in question here; their usefulness as direct
autobiography is. Modern criticism disputes their direct authorship in
sharply different degrees---least for 2~Thessalonians and Colossians, where
substantial scholarly minorities (in some quarters majorities) still defend
Pauline composition; more for Ephesians; most for the Pastorals. The
scale of the discussion can be stated from a checked witness: teaching on
the Pastorals during the Pauline Year, Benedict~XVI said plainly that
``most exegetes today are of the opinion that these Letters would not have
been written by Paul himself, but would have come from the `Pauline
School','' while noting that parts of 2~Timothy ``appear so authentic that
they could have come only from the heart and mouth of the Apostle''
(General Audience, 28 January 2009). This study therefore uses the six in a
bounded way: as first-rank evidence for how Paul was remembered, extended,
and institutionalized within a generation or two of his death; as possible
but not certain evidence for events after Acts ends (2~Timothy's Roman
imprisonment, Titus's Crete and Nicopolis); and never as the sole support
for a biographical claim. Hebrews, anonymous and stylistically remote, was
attached to Paul's name in parts of the ancient church but makes no
autobiographical claim and is used here only in the history of his
reception.

\witnesslimit{Ancient authorship judgments are also part of the record: the
Muratorian fragment (Rome?, conventionally c.~170--200) already counts
thirteen letters, ``to no more than seven churches by name'' plus Philemon,
Titus, and two to Timothy, and already knows of forgeries---letters ``to
the Laodiceans, and another to the Alexandrians, forged under the name of
Paul.'' Letters in Paul's name that Paul did not write were a live
first--second-century phenomenon; 2~Thessalonians itself warns against being shaken ``by epistle, as sent
from us'' (2~Thess 2:2) and closes with an authenticating signature, ``the
salutation of Paul with my own hand; which is the sign in every epistle''
(2~Thess 3:17).}

\subsection{Acts as a second witness}

The Acts of the Apostles devotes more than half its length to Paul
(chiefly Acts 7:58--8:3; 9; 11:25--30; 13--28). It is a first-century
canonical narrative, ancient tradition's ``Luke, the most dear physician''
(Col 4:14) writing a generation after the events, and from Acts 16 onward
it intermittently narrates in the first person plural (``immediately we
sought to go into Macedonia,'' Acts 16:10)---passages read since Irenaeus
as the diary of a companion: Luke was ``inseparable from Paul, and his
fellow-labourer in the Gospel'' (\work{Against Heresies} 3.14.1). Whatever
the ``we''-source is, Acts as a whole is a designed theological history:
its speeches are compositions in character, not transcripts; its journeys
are shaped into three great circuits; its Paul works miracles, keeps the
law, and converses urbanely with proconsuls and kings. Strikingly, Acts
never once mentions that Paul wrote letters---the surest sign that it is
an independent witness, not a narrative spun from the epistles. This study
treats Acts as class-N evidence: near-contemporary, external, indispensable
for the itinerary and for data the letters never state (Tarsus, Gamaliel,
Roman citizenship, the name Saul), and always to be checked against the
letters where they overlap.

\subsection{Where Acts and the letters pull apart}

Four tensions matter enough to govern the whole reconstruction, and each
is examined in place below. (1)~\textbf{The Jerusalem visits.} Paul swears
--- ``before God, I lie not'' (Gal 1:20)---that after the revelation he did
not go up to Jerusalem for three years, then saw only Cephas and James, and
came again only ``after fourteen years'' (Gal 1:18--19; 2:1). Acts has the
new convert brought to ``the apostles'' by Barnabas soon after Damascus
(9:27), then a famine-relief visit (11:29--30), then the council visit
(15). Whether Galatians~2 corresponds to Acts~15 or to Acts~11, and what
becomes of the extra visit, is a genuine, unresolved critical problem, not
a detail. (2)~\textbf{The Damascus escape.} Paul: ``At Damascus, the
governor of the nation under Aretas the king, guarded the city\ldots\ and
through a window in a basket was I let down by the wall'' (2~Cor
11:32--33). Acts: ``the Jews consulted together to kill him'' and watched
the gates (9:23--24). Same basket, same wall, different pursuers---a
textbook case of two tellings of one memory, each shaped by its teller's
frame. (3)~\textbf{Emphasis and omission.} Acts' Paul is a triumphant
wonder-worker; the letters' Paul catalogues beatings, shipwrecks, and ``a
sting of my flesh'' (2~Cor 11:24--27; 12:7), and his opponents said ``his
bodily presence is weak'' (2~Cor 10:10). Acts omits the collection's
delivery, the crisis with Corinth, the conflict at Antioch, and the letters
themselves. (4)~\textbf{Title and status.} The letters open with, argue
from, and die on the title \emph{apostle}---``Am not I an apostle?''
(1~Cor 9:1)---which Paul grounds directly in having seen the risen Lord;
Acts, which reserves the word almost entirely for the Twelve, applies it
to Paul only twice in passing (``the apostles Barnabas and Paul,'' Acts
14:4, 13 DR numbering), while conversely the Roman citizenship on which
Acts' whole legal drama turns is never once mentioned in the letters.
Each witness, that is, has its own theory of who Paul was. Where the two
witnesses conflict, the participant outranks the
narrator; where they merely differ in selection, both are kept, severally.

\subsection{The rest of the record}

After the New Testament the witnesses stratify cleanly, and the strata are
used strictly by date and genre. Early received tradition (class E):
1~Clement, written from Rome conventionally c.~95--96, within one
generation of the death; Ignatius of Antioch (c.~107--110 traditional) and
Polycarp of Smyrna (c.~110--130s), who both invoke Paul by name; the
second-century reports preserved at exact loci in Eusebius---Dionysius of
Corinth (c.~170), the Roman presbyter Gaius (c.~200) on the ``trophies'' of
the apostles, Origen (c.~230) on the manner of death; Irenaeus and
Tertullian around 200. Later hagiography and legend (class L): the
apocryphal \work{Acts of Paul} (c.~160 by Tertullian's notice), source of
Thecla, of the famous physical description, and of the earliest martyrdom
narrative; the fourth--fifth-century martyrdom romances that supplied
Lucina's estate and the three fountains; the forged Paul--Seneca
correspondence that Jerome read. Official and archaeological reception
(classes M and K): the basilica of St.~Paul Outside the Walls and its
documented tomb presentation; the 2002--2006 excavation of the sarcophagus
and Benedict~XVI's announcement of its scientific probing (28 June 2009);
the liturgical calendar. Modern critical reconstruction (class K) is cited
as named positions---above all for the chronology, where a single
epigraphic control, the Gallio inscription from Delphi, published in the
modern era and studied classically by Adolf Deissmann, supplies the one
absolute date in Paul's life. Project synthesis (class S) is the
responsibility of this study alone and is never attributed to a source.
